\chapter[Presentación]{Presentación}
\label{cp:presentacion}

\insertminitoc
\parindent0pt

\section{Planteamiento del Problema}

Según la Organización Mundial de la Salud (\acrshort{oms}), más de 1\,500 millones de personas en
el mundo presentan algún grado de pérdida auditiva, de las cuales aproximadamente 430
millones requieren atención especializada. De este total, se estima que al menos 34 millones
son niños \cite{who_deafness_2024}. Para esta población, la lengua de señas constituye el principal medio de
comunicación y expresión, configurando una forma de interacción con valores y costumbres
propias que difieren de las de la población oyente.

En Honduras, la comunidad sorda está organizada en torno a la Asociación de Sordos de
Honduras (\acrshort{ash}), institución que ha impulsado el desarrollo y la documentación del
Lenguaje de Señas Hondureño (\acrshort{lesho}), el cual fue reconocido oficialmente como la lengua
natural de la comunidad sorda mediante el Decreto Legislativo No. 321-2013 \cite{congreso_decreto_2013}, \cite{ash_diccionario_2019}. Sin
embargo, según datos de la Comisión Nacional de los Derechos Humanos (\acrshort{conadeh}),
en Honduras existen aproximadamente 100\,000 personas con discapacidad auditiva, de las
cuales únicamente alrededor del 1\,\% cuenta con un carné oficial de identificación, lo que
evidencia el bajo nivel de reconocimiento institucional de esta población \cite{conadeh_discriminados}.

La brecha de comunicación entre niños sordos y personas oyentes se manifiesta en
prácticamente todos los entornos cotidianos, incluyendo el hogar, la comunidad y los
espacios de atención en salud. La mayoría de los niños sordos nacen en hogares de padres
oyentes que no conocen el \acrshort{lesho}, lo que retrasa el acceso temprano a una lengua completa
y limita su desarrollo socioemocional \cite{williams_encuesta_2015}. A esto se suma la escasez de intérpretes
certificados fuera de entornos educativos especializados, especialmente en zonas rurales
o de bajos recursos donde su presencia es prácticamente inexistente.

A pesar de que existen soluciones tecnológicas para el reconocimiento de otras lenguas de
señas, como la Lengua de Señas Americana o la Lengua de Señas Brasileña, no existe a la
fecha una aplicación móvil que permita a un niño sordo expresarse en \acrshort{lesho} hacia una
persona oyente en tiempo real, ni que facilite la comunicación en la dirección inversa
representando el español de forma visual en la lengua materna del niño. Las herramientas
comerciales disponibles no contemplan el vocabulario ni las particularidades del \acrshort{lesho}, y
su idioma de interfaz y su costo representan barreras adicionales para la población
hondureña \cite{sadik_survey_2025}.

\subsection{Preguntas de Investigación}

¿Es posible desarrollar una aplicación móvil que reconozca gestos del Lenguaje de Señas
Hondureño en tiempo real utilizando visión por computadora y aprendizaje profundo
ejecutado directamente en el dispositivo?

¿Puede una solución bidireccional facilitar la comunicación entre niños sordos y personas
oyentes sin requerir conocimiento previo del \acrshort{lesho} por parte de ninguna de las dos partes?

¿Qué nivel de precisión y usabilidad puede alcanzar un sistema de este tipo desarrollado
con un conjunto de datos propio de \acrshort{lesho} de alcance acotado?

\subsection{Justificación del Estudio}

La igualdad en la comunicación es un derecho que debe garantizarse para todas las
personas. Los niños sordos en Honduras enfrentan barreras cotidianas que van más allá
del ámbito educativo, afectando su capacidad de expresarse y comprender en entornos
familiares, comunitarios y de salud. La dependencia exclusiva de intérpretes humanos
profundiza esta desigualdad en contextos donde su presencia es escasa o inexistente \cite{congreso_decreto_2013}.

La relevancia de este trabajo radica en explorar una solución tecnológica contextualizada
a la realidad hondureña, dado que la mayoría de herramientas existentes han sido diseñadas
para otras lenguas de señas y otros entornos. El desarrollo de una aplicación móvil
bidireccional que opere sin conexión a internet es especialmente pertinente en un país
con cobertura de datos móviles desigual, donde no puede garantizarse el acceso a servicios
en la nube de forma permanente.

El estudio tiene impacto tanto académico como social. Desde el punto de vista académico,
la construcción de un conjunto de datos de señas del \acrshort{lesho} orientado al entrenamiento de modelos
de inteligencia artificial constituye un aporte al campo tecnológico local, ya que durante
la revisión bibliográfica realizada no se encontraron antecedentes documentados de un
recurso de este tipo para esta lengua.
Desde el punto de vista social, la aplicación busca reducir la brecha de comunicación entre
niños sordos y personas oyentes, aportando una herramienta de inclusión comunicativa
de bajo costo y sin dependencia de infraestructura externa.

\subsection{Viabilidad del Estudio}

La viabilidad del estudio se analiza en función de los siguientes puntos.

\begin{enumerate}
    \item Disponibilidad de recursos.
    \begin{itemize}
        \item \textbf{Humanos:} la investigación será desarrollada por el tesista con la guía del asesor
        correspondiente, y contará con la colaboración de personas con conocimiento de
        \acrshort{lesho} vinculadas a fundaciones de atención a personas con discapacidad en el
        departamento de Comayagua.
        \item \textbf{Tecnológicos:} se dispone de herramientas de código abierto para el desarrollo,
        incluyendo Flutter, MediaPipe y TensorFlow Lite, las cuales no requieren licencias
        de pago y cuentan con documentación oficial extensa \cite{google_mediapipe_2024}, \cite{google_tflite_2024}, \cite{google_flutter_2024}.
        \item \textbf{Datos:} el conjunto de datos de señas será construido durante la investigación siguiendo un
        protocolo de captura documentado, con referencia al diccionario de \acrshort{lesho}
        publicado por la \acrshort{ash} \cite{ash_diccionario_2019}.
    \end{itemize}

    \item Alcance del estudio.

    La investigación se centrará en desarrollar una aplicación móvil funcional que
    reconozca el alfabeto dactilológico y un vocabulario de 50 señas dinámicas, y que
    permita representar visualmente señas del \acrshort{lesho} a partir de texto en español. El
    estudio no pretende abarcar la totalidad del vocabulario del \acrshort{lesho} ni contempla
    aspectos gramaticales como expresiones faciales o lenguaje corporal. La evaluación
    se realizará en condiciones controladas con usuarios con conocimiento de la lengua.

    \item Implicaciones del estudio.
    \begin{itemize}
        \item \textbf{Académicas:} contribuye al conocimiento en el campo del reconocimiento de gestos
        y la accesibilidad tecnológica aplicada a la comunidad sorda hondureña.
        \item \textbf{Tecnológicas:} promueve el desarrollo de soluciones con inteligencia artificial
        ejecutada en el dispositivo, sin dependencia de servidores externos.
        \item \textbf{Sociales:} fomenta la inclusión comunicativa de niños sordos en entornos
        cotidianos, aportando una herramienta accesible y de bajo costo.
    \end{itemize}

    \item Consecuencias del estudio.

    Los resultados obtenidos pueden servir como base para investigaciones futuras que
    amplíen el vocabulario reconocido, incorporen componentes no manuales del \acrshort{lesho}
    o extiendan la solución a otros contextos de uso. La metodología de construcción
    del conjunto de datos y el proceso de entrenamiento de modelos quedan documentados de
    forma reproducible para facilitar su extensión.
\end{enumerate}

\section{Objetivo General}

Desarrollar una aplicación de inteligencia artificial para la comunicación entre niños
sordos y oyentes mediante el reconocimiento de señas \acrshort{lesho}, ejecutada de forma
completamente local en el dispositivo móvil.

\section{Objetivos Específicos}

\begin{enumerate}
    \item Recopilar un conjunto de datos de señas \acrshort{lesho} que incluya el alfabeto dactilológico
    completo y un vocabulario de señas dinámicas de uso cotidiano, siguiendo un protocolo
    de captura documentado y reproducible.

    \item Diseñar e implementar modelos de clasificación de señas estáticas y dinámicas
    optimizados para su ejecución en dispositivos móviles con recursos computacionales
    limitados.

    \item Desarrollar una aplicación que integre el reconocimiento automático de señas y un
    diccionario visual de señas pregrabadas, operando sin dependencia de servicios en la
    nube.

    \item Evaluar el desempeño técnico del sistema en términos de exactitud de clasificación y
    latencia de inferencia, y valorar la usabilidad de la aplicación mediante pruebas con
    usuarios que incluyan personas con conocimiento de \acrshort{lesho}.
\end{enumerate}
